\documentclass[10pt,a4paper]{amsart}
\usepackage[margin=19mm]{geometry}
\usepackage{amsmath,amssymb,mathtools}
\usepackage[scr=rsfs,cal=euler]{mathalfa}
\usepackage{xfrac}
\usepackage[unicode,hidelinks,bookmarks=false]{hyperref}
\newcommand{\ie}{i.e.\ }
\newcommand{\cf}{cf.\ }
\newcommand{\resp}{respectively\ }
\newcommand{\Subsection}{\subsection}
\providecommand{\nobreakdash}{\nobreak\hbox{-}}
\setlength{\emergencystretch}{2em}
\multlinegap0pt
\usepackage{bm}
\usepackage{bbm}
\usepackage{dsfont}
\usepackage{xspace}
\usepackage{cancel}
\usepackage{mathtools}
\usepackage{amsthm}
\makeatletter
\@ifundefined{thm}{\newtheorem{thm}{Thm}}{}
\@ifundefined{prop}{\newtheorem{prop}[thm]{Proposition}}{}
\makeatother
\newcommand{\N}{\mathbb{N}}
\renewcommand{\geq}{\geqslant}
\renewcommand{\leq}{\leqslant}
\title[Density estimates for a (non)local variational model with de]{Density estimates for a (non)local variational model with degenerate double-well potential}
\author{Serena Dipierro, Alberto Farina, Giovanni Giacomin, Enrico Valdinoci}
\date{Source: arXiv:2506.22193v1 (CC BY 4.0)}
\begin{document}
\maketitle

\noindent\emph{Source and licence.} Density estimates for a (non)local variational model with degenerate double-well potential. Version arXiv:2506.22193v1 (CC BY 4.0), distributed under the Creative Commons Attribution 4.0 International licence, as recorded in the arXiv licence metadata for this exact version and shown on the abstract page https://arxiv.org/abs/2506.22193v1. Full rights notice: licenses/arxiv-2506.22193-CC-BY-4.0.txt.\par\smallskip

\noindent\textit{Editorial scope.} Counted unit: the Prop whose source label is \texttt{colllise} in the article (source0.tex lines 2248--2259, source label \texttt{colllise}), delivered together with the article's own complete original proof of it (source0.tex lines 2262--2358, one \texttt{proof} environment of 97 lines). No mathematics is written, repaired or re-derived: statement and proof are the article's own text line for line. Required context retained uncounted: potntiald2 (lines 152--154). Every definition, display and label of those lines is retained unchanged. Omitted from this package and not redistributed: 1--151, 155--2247, 2260--2261, 2359--3758. Nothing inside a retained excerpt is omitted or altered: every retained body is delivered exactly as the article's lines are written, so no omission receipt of any kind accompanies this package. Citations: the retained excerpts carry no citation command at all, so no bibliography entry of the article is transcribed and no reference is redistributed. Reference adaptation: none was needed and none was made; every reference inside the delivered unit resolves inside it, so no unresolved reference or citation is produced. Printed slips kept literal: no printed word, symbol, hypothesis or number of the counted statement or of its proof was altered. Layout and class: the article's own class is \texttt{article} with packages including geometry, amssymb, amsmath, amsthm, graphicx, epstopdf, color, esint, amsrefs, inputenc, miama, fontenc, mathtools, authblk; the wrapper is the lane's pinned \texttt{amsart} editorial wrapper at 10pt on A4, which restates the theorem-like environments the retained text needs and adds the article's own macro definitions that the retained text needs (\texttt{\textbackslash N}, \texttt{\textbackslash geq}, \texttt{\textbackslash leq}), with extra wrapper packages (bm, bbm, dsfont, xspace, cancel, mathtools). The delivered numbering is the unit's own. Assets: no retained span contains a figure environment, a graphics-inclusion command, a PDF-page inclusion, a table or a TikZ picture; no image file is copied into this package, the package declares no asset dependency and no image file is redistributed with it. Licence: version arXiv:2506.22193v1 of this article is distributed under the Creative Commons Attribution 4.0 International licence, as recorded in the arXiv licence metadata for this exact version and shown on the abstract page https://arxiv.org/abs/2506.22193v1. A full rights notice is delivered at licenses/arxiv-2506.22193-CC-BY-4.0.txt. Characters: every retained excerpt is pure ASCII, so the delivered engine, XeLaTeX with its default Latin Modern text font, reports no missing character.
\begin{equation}\label{potntiald2}
W(t)-W(r)\geq c_1(t-r)\left|1+r\right|^{m-1}+c_1(t-r)^k,
\end{equation}

\begin{prop}\label{colllise}
Let $m\in (1,+\infty)$ and $W_m(x):=(1-x^2)^m$ with $x\in[-1,1]$. Then, for every $k\in \left\lbrace 1,\dots, [m] \right\rbrace$,
\begin{equation}\label{telenious}
W_m^{(k)}(x)=W_{m-k}(x)P_m^k(x).
\end{equation}

Also, for every $m\in (1,+\infty)$ and $c\in (0,1)$ there exists~$\hat{q}_{m,c}\in (0,1)$ such that, for every $k\in \left\lbrace 1,\dots, \left[m\right] \right\rbrace$ and $x\in [-1,-1+\hat{q}_{m,c}]$, 
\begin{equation}\label{term-0987}
P_m^k(x)\geq c.
\end{equation}
In particular, $W_m(x)$ satisfies~\eqref{potntiald2} for some $c_1\in (0,+\infty)$ and $q\in (0,1)$. 
\end{prop}

\begin{proof} 
We notice that, for every $m\in (1,+\infty)$,
\begin{equation}\label{klkp.mtv5}
W_m'(x)=-2xm(1-x^2)^{m-1}=P_m^1(x)W_{m-1}(x).
\end{equation}

Let us now proceed recursively and assume that $m\in [2,+\infty)$ and that~\eqref{telenious} holds for every $j\leq k-1<\left[m\right]$. Then, 
\begin{equation*}
\begin{split}
W_m^{(k)}(x)&=\left(W_m^{(k-1)}(x)\right)'\\
&=\left(W_{m-k+1}(x)P_m^{k-1}(x)\right)'\\
&=W_{m-k+1}'(x)P_m^{k-1}(x)+W_{m-k+1}(x){P_m^{k-1}}'(x)\\
&=P_{m-k+1}^1(x)W_{m-k}(x)P_m^{k-1}(x)+W_{m-k+1}(x){P_m^{k-1}}'(x)\\
&=W_{m-k}(x)\left(P_{m-k+1}^1(x)P_m^{k-1}(x)+(1-x^2){P_m^{k-1}}'(x)\right)\\
&=W_{m-k}(x)P_m^k(x).
\end{split}
\end{equation*}
This concludes the proof of~\eqref{telenious}.

Now we prove~\eqref{term-0987}. For this, we consider $0<q_{\left[m\right]}\leq \dots\leq q_1<1$ which will be chosen as follows. 

For any $m\in (1,+\infty)$, we choose $q_1:=1-\frac{c}{2m}$, from which it follows that, for every $x\in \left[-1,-1+q_1\right]$,
\begin{equation*}
P_m^1(x)=-2xm\geq -2m(-1+q_1)=2m(1-q_1)=c.
\end{equation*}
Now we assume that $P_m^j(x)\geq c$ for every $j\leq k-1<\left[m\right]$ and $x\in \left[-1,-1+q_j\right]$. We need to show that there exists~$q_k\in\left(0, q_{k-1}\right)$ such that $P_m^k(x)\geq c$ for every $x\in [-1,-1+q_k]$. In order to do so, we let 
\begin{equation}
q_k:=\min \left\lbrace q_{k-1}, \frac{2c(m-k+1)-c}{2c(m-k+1)+2\left\|{P_m^{k-1}}'\right\|_{L^\infty([-1,0])}} \right\rbrace
\end{equation}
and, for every $x\in [-1,-1+q_k]\subset [-1,-1+q_{k-1}]$,
\begin{equation*}
\begin{split}
P_{m}^k(x)&=P_{m-k+1}^1(x)P_m^{k-1}(x) +(1-x^2){P_m^{k-1}}'(x)\\
&\geq P_{m-k+1}^1(x)P_m^{k-1}(x)-(1-x^2)\left\|{P_m^{k-1}}'\right\|_{L^\infty([-1,0])}\\
& \geq 2c(1-q_k)(m-k+1)-2q_k\left\|{P_m^{k-1}}'\right\|_{L^\infty([-1,0])}\\
&\geq c.
\end{split}
\end{equation*}
In particular, by choosing $\hat{q}_{c,m}:=q_{\left[m\right]}$ we establish the validity of~\eqref{term-0987}, as desired.

Thus, as a consequence of~\eqref{telenious} and~\eqref{term-0987},
we find that 
\begin{equation*}
\begin{split}
W_m(t)-W_m(r)&=\int_{r}^tW_m'(\xi)\,d\xi\\
&=\sum_{j=1}^{k-1} \frac{W_m^{(j)}(r)}{j!}(t-r)^j+\int_{r}^t\dots \int_{r}^l W_m^{(k)}(z)\,dz\dots\,d\xi\\
&\geq W_m^{(1)}(r)(t-r)  + \frac{c}{k!}(t-r)^k\\
&=(1-r^2)^{m-1}P_1^m(r)(t-r)+\frac{c}{k!}(t-r)^k\\
&\geq c(1-r)^{m-1}(1+r)^{m-1}(t-r)+\frac{c}{k!}(t-r)^k\\
&\geq c(2-\hat{q}_{c,m})^{m-1}(1+r)^{m-1}(t-r)+\frac{c}{k!}(t-r)^k,
\end{split}
\end{equation*}
as long as $m\in \N$, $c\in (0,1)$, $-1\leq r\leq t\leq -1+\hat{q}_{c,m}$ and $k=m$.

This shows that $W_m(x)$ satisfies~\eqref{potntiald2} for every $m\in\N$  with $q:=\hat{q}_{c,m}$ and 
\begin{equation}\label{c_1}
c_1:=\min\left\lbrace  c(2-q)^{m-1} ,\frac{c}{k!}    \right\rbrace.
\end{equation}

Additionally, if $m\in (1,+\infty)\setminus \N$ and $c\in (0,1)$ we define $q_1\in (0,1)$ as the unique solution to 
\begin{equation*}
2c(1-q_1)(m-\left[m\right])(2 q_1)^{m-\left[m\right]-1} =c+ (2 q_1)^{m-\left[m\right]}\left\|{P_m^{\left[m\right]}}'\right\|_{L^\infty([-1,0])}.
\end{equation*}
We observe that, for every $t\in (0,q_1)$, 
\begin{equation}\label{jvecogdv4}
2c(1-t)(m-\left[m\right])(2 t)^{m-\left[m\right]-1}>c+ (2 t)^{m-\left[m\right]}\left\|{P_m^{\left[m\right]}}'\right\|_{L^\infty([-1,0])}.
\end{equation}
Then, we set  
\begin{equation}\label{eq!}
q:=\min \left\lbrace \hat{q}_{c,m}, q_1  \right\rbrace
\end{equation}
and we deduce from~\eqref{jvecogdv4},~\eqref{telenious} and~\eqref{term-0987} that,
when $k=\left[m\right]+1$,
\begin{equation*}
\begin{split}
W_m^{(k)}(x)&=\left(W_m^{([m])}(x)\right)'\\
&=\left((1-x^2)^{m-\left[m\right]}P_m^{\left[m\right]}(x)\right)'\\
&=-2x(m-\left[m\right])\left(1-x^2\right)^{m-\left[m\right]-1}P_m^{\left[m\right]}(x)+ (1-x^2)^{m-\left[m\right]}{P_m^{\left[m\right]}}'(x)\\
&\geq 2c(1-\tilde{q})(m-\left[m\right])(2\tilde{q})^{m-\left[m\right]-1}-(2\tilde{q})^{m-\left[m\right]}\left\|{P_m^{\left[m\right]}}'\right\|_{L^\infty([-1,0])}\\
&\geq c,
\end{split}
\end{equation*}
for every $x\in [-1,-1+q]$

Because of this, we have that, for every $-1\leq r\leq t\leq -1+q$,
\begin{equation*}
\begin{split}
W_m(t)-W_m(r)&=\int_{r}^tW_m'(\xi)\,d\xi\\
&=\sum_{j=1}^{k-1} \frac{W_m^{(j)}(r)}{j!}(t-r)^j+\int_{r}^t\dots \int_{r}^l W_m^{(k)}(z)\,dz\dots\,d\xi\\
&\geq W_m^{(1)}(r)(t-r)  + \frac{c}{k!}(t-r)^k\\
&=(1-r^2)^{m-1}P_1^m(r)(t-r) \frac{c}{k!}(t-r)^k\\
&\geq c(1-r)^{m-1}(1+r)^{m-1}(t-r)+\frac{c}{k!}(t-r)^k\\
&\geq c(2-q)^{m-1}(1+r)^{m-1}(t-r)+\frac{c}{k!}(t-r)^k.
\end{split}
\end{equation*}
This proves that $W_{m}$ satisfies~\eqref{potntiald2} for every $m\in (1,+\infty)\setminus \N$ with $c_1$ and $q$ respectively as in~\eqref{c_1} and~\eqref{eq!}.  
\end{proof}

\end{document}
